\Needspace{14\baselineskip}
\fancyhead[L]{\small ZWEITER KOMET VON 1805.}
\begin{center}
ZWEITER KOMET DES JAHRES 1805.

\rule{0.6\linewidth}{0.4pt}

Monatliche Correspondenz. Band XIII. S. 83–91. Januar 1806.

\rule{0.6\linewidth}{0.4pt}

—\quad—\quad—
\end{center}

Dr. \textsc{Gauss} beobachtete diesen Kometen am 8. Dezember. Er wurde mit 359 im Wassermann und einem anderen Stern aus der \emph{Hist. célest.} verglichen.

\begin{center}
\begingroup\fontsize{9}{11}\selectfont\setlength{\tabcolsep}{3pt}\renewcommand{\arraystretch}{1.1}
\begin{tabular}{r|r|r|r}
1805 & \shortstack{Mittlere Zeit in\\Braunschweig} & \shortstack{Scheinbare\\Rektaszension} & \shortstack{Scheinbare südliche\\Deklination}\\ \hline
8. Dezember & $6^{\mathrm u}\ 55^{\mathrm m}\ 17^{\mathrm s}$ & $353^\circ\ 7'\ 40''$ & $23^\circ\ 36'\ 24''$ \\
8 & $7\ 58\ 3$ & $352\ 57\ 57$ & $24\ 6\ .\ .$ \\
\end{tabular}\endgroup
\end{center}

\par\indent Die zweite Deklination ist nur eine Schätzung. Aus dieser Beobachtung, der ersten Beobachtung von \textsc{Bouvard} vom 16. November und der oben angeführten Beobachtung von Dr. \textsc{Olbers} vom 2. Dezember berechnete Dr. \textsc{Gauss} die folgenden parabolischen Elemente.

\begin{center}
\begingroup\fontsize{8}{10}\selectfont\setlength{\tabcolsep}{3pt}\renewcommand{\arraystretch}{1.1}
\begin{tabular}{lr}
Durchgang durch das Perihel am 31. Dezember 1805 , & $7^{\mathrm u}\ 20^{\mathrm m}\ 39^{\mathrm s}$ Seeberger Zeit \\
Länge des Perihels , . . . . . . . . . . . . & $109^\circ\ 23'\ 40''$ \\
Länge des aufsteigenden Knotens . . . . . . . . . . & $250\ 33\ 14$ \\
Neigung der Bahn . , . . . . . \raisebox{.3ex}{.} . . . . . . . & $16\ 33\ 33$ \\
Logarithmus des kleinsten Abstands . . . . . . . . & $9.9502477$ \\
Bewegung . . . . . . . . . . . . . . . . . . & rechtläufig. \\
\end{tabular}\endgroup
\end{center}

\Needspace{24\baselineskip}
\par\indent Nach der Methode aus dem Aufsatz im Maiheft der M. C. von 1804 [S. \hyperlink{native-99}{94} und \hyperlink{native-125}{120} dieses Bandes] werden die Koordinaten durch die folgenden Formeln dargestellt.

\[
\begin{gathered}
x = \frac{\alpha \sin(v + 198^\circ\,37'\,20'')}{\cos\frac{1}{2}v^2},\\[4pt]
y = \frac{\ObservedBeta \sin(v + 103^\circ\,41'\,12'')}{\cos\frac{1}{2}v^2},\\[4pt]
z = \frac{\gamma \sin(v + 148^\circ\,41'\,54'')}{\cos\frac{1}{2}v^2}
\end{gathered}
\]

\noindent wobei $v$ die Anomalie des Kometen bezeichnet und Log. $\alpha = 9.933968$, Log. $\ObservedBeta = 9.930525$, Log. $\gamma = 9.551021$.

Aus seinen Elementen bestimmt Dr. \textsc{Gauss} unter Berücksichtigung der Parallaxe für den Abstand des Kometen von der Erde am 16. Nov. 0.24990, am 2. Dez. 0.08913, am 8. Dez. bei der I. Beob. 0.04850, am 8. Dez. bei der II. Beob. 0.04841.

Für die Lichtstärke ergibt sich dagegen am 16. November 16.6, am 2. Dezember 117.8 und am 8. Dezember 439.2, wenn man die Lichtstärke im Abstand 1 von der Erde und von der Sonne als Einheit annimmt.

\begin{center}
—\quad—\quad—
\end{center}

\hfill[\textsc{von Zach.}]
